\textbf{Phần 1. Trắc nghiệm lựa chọn}
\setcounterpageref{ex}{0}

\begin{ex} %[3X15N73]
Khi kiểm đếm số lượng học sinh thích các môn thể thao, ta thực hiện công việc
\choice
{Vẽ tranh minh họa}
{Tính chu vi sân bóng}
{\True Thu thập và ghi chép số liệu}
{Làm tròn số đến hàng chục}
\loigiai{
Đó là hoạt động thu thập và ghi chép số liệu.\\
\textbf{Đáp án của câu hỏi là C}
}
\end{ex}

\begin{ex} %[3X15N73]
Cho bảng số liệu số cây trồng được của 3 lớp: Lớp 3A: 45 cây; Lớp 3B: 50 cây; Lớp 3C: 40 cây. Lớp trồng được nhiều cây nhất là
\choice
{Lớp 3A}
{Lớp 3C}
{Ba lớp bằng nhau}
{\True Lớp 3B}
\loigiai{
Lớp 3B trồng được nhiều nhất (50 cây).\\
\textbf{Đáp án của câu hỏi là D}
}
\end{ex}

\begin{ex} %[3X15H73]
Dựa vào bảng số liệu trên, tổng số cây cả 3 lớp trồng được là
\choice
{$130$ cây}
{\True $135$ cây}
{$140$ cây}
{$145$ cây}
\loigiai{
$45 + 50 + 40 = 135$ cây.\\
\textbf{Đáp án của câu hỏi là B}
}
\end{ex}

\begin{ex} %[3X15H73]
Lớp 3B trồng nhiều hơn lớp 3C số cây là
\choice
{\True $10$ cây}
{$5$ cây}
{$15$ cây}
{$8$ cây}
\loigiai{
$50 - 40 = 10$ cây.\\
\textbf{Đáp án của câu hỏi là A}
}
\end{ex}

\begin{ex} %[3X15H73]
Một cửa hàng bán hoa thống kê số hoa bán trong ngày: Hoa hồng: 120 bông, Hoa cúc: 85 bông, Hoa ly: 60 bông. Loại hoa nào bán được ít nhất?
\choice
{Hoa cúc}
{\True Hoa ly}
{Hoa hồng}
{Cả 3 loại bằng nhau}
\loigiai{
Hoa ly bán được ít nhất (60 bông).\\
\textbf{Đáp án của câu hỏi là B}
}
\end{ex}

\begin{ex} %[3X15H73]
Số hoa hồng bán nhiều hơn hoa cúc là
\choice
{$30$ bông}
{$40$ bông}
{$25$ bông}
{\True $35$ bông}
\loigiai{
$120 - 85 = 35$ bông.\\
\textbf{Đáp án của câu hỏi là D}
}
\end{ex}

\begin{ex} %[3X15V73]
Số học sinh giỏi các khối của một trường: Khối 1: 80 em; Khối 2: 95 em; Khối 3: 110 em; Khối 4: 105 em; Khối 5: 120 em. Khối có số học sinh giỏi nhiều nhất hơn khối ít nhất là
\choice
{$35$ em}
{$45$ em}
{\True $40$ em}
{$50$ em}
\loigiai{
Khối 5 nhiều nhất (120 em), Khối 1 ít nhất (80 em). Hiệu: $120 - 80 = 40$ em.\\
\textbf{Đáp án của câu hỏi là C}
}
\end{ex}

\begin{ex} %[3X15V73]
Một thư viện thống kê số sách cho mượn trong 4 ngày: Thứ Hai: 45 cuốn; Thứ Ba: 60 cuốn; Thứ Tư: 55 cuốn; Thứ Năm: 70 cuốn. Tổng số sách đã cho mượn trong 4 ngày là
\choice
{\True $230$ cuốn}
{$220$ cuốn}
{$240$ cuốn}
{$250$ cuốn}
\loigiai{
$45 + 60 + 55 + 70 = 230$ cuốn.\\
\textbf{Đáp án của câu hỏi là A}
}
\end{ex}

\begin{ex} %[3X15V73]
Số sách mượn ngày thứ Năm nhiều hơn ngày thứ Hai là
\choice
{$20$ cuốn}
{$30$ cuốn}
{$15$ cuốn}
{\True $25$ cuốn}
\loigiai{
$70 - 45 = 25$ cuốn.\\
\textbf{Đáp án của câu hỏi là D}
}
\end{ex}

\begin{ex} %[3X15V73]
Thống kê điểm kiểm tra Toán của tổ 1: có 3 bạn được 10, 4 bạn được 9, 2 bạn được 8 và 1 bạn được 7. Tổ 1 có tất cả bao nhiêu bạn?
\choice
{\True $10$ bạn}
{$9$ bạn}
{$11$ bạn}
{$8$ bạn}
\loigiai{
$3 + 4 + 2 + 1 = 10$ bạn.\\
\textbf{Đáp án của câu hỏi là A}
}
\end{ex}

\begin{ex} %[3X15V73]
Số bạn đạt điểm từ 9 trở lên trong tổ 1 là
\choice
{$6$ bạn}
{$8$ bạn}
{\True $7$ bạn}
{$9$ bạn}
\loigiai{
Số bạn đạt điểm 9 hoặc 10 là: $3 + 4 = 7$ bạn.\\
\textbf{Đáp án của câu hỏi là C}
}
\end{ex}

\begin{ex} %[3X15C73]
Thống kê số xe đạp bán được trong 3 tháng đầu năm: Tháng 1: 150 xe, Tháng 2: 120 xe, Tháng 3 bán được nhiều hơn tháng 2 là 40 xe. Cả 3 tháng bán được số xe là
\choice
{$420$ xe}
{\True $430$ xe}
{$440$ xe}
{$450$ xe}
\loigiai{
Tháng 3 bán: $120 + 40 = 160$ xe.\\
Cả 3 tháng: $150 + 120 + 160 = 430$ xe.\\
\textbf{Đáp án của câu hỏi là B}
}
\end{ex}

\textbf{Phần 2. Trắc nghiệm đúng sai}
\setcounterpageref{ex}{0}

\begin{ex} %[3X15H73]
Bảng số liệu số con vật nuôi trong một trang trại: Gà: 350 con, Vịt: 280 con, Heo: 120 con, Bò: 50 con. Xét tính đúng sai:
\choiceTF
{\True Gà là con vật được nuôi nhiều nhất (350 con).}
{\True Bò là con vật được nuôi ít nhất (50 con).}
{Tổng số gà và vịt là 600 con.}
{\True Tổng số gà và vịt là 630 con ($350 + 280 = 630$).}
\loigiai{
\begin{itemize}
    \item Gà nhiều nhất: Đúng.
    \item Bò ít nhất: Đúng.
    \item $350 + 280 = 630 \neq 600$. Ý 3 Sai, Ý 4 Đúng.
\end{itemize}
\begin{itemize}
    \item a) \textbf{Đúng}.
    \item b) \textbf{Đúng}.
    \item c) \textbf{Sai}.
    \item d) \textbf{Đúng}.
\end{itemize}
}
\end{ex}

\begin{ex} %[3X15V73]
Thống kê số lượng cây xanh trồng trong vườn trường: Cây bàng: 15 cây, Cây phượng: 12 cây, Cây bằng lăng: 8 cây. Xét tính đúng sai:
\choiceTF
{\True Vườn trường trồng tất cả 35 cây xanh.}
{\True Cây bàng nhiều hơn cây bằng lăng là 7 cây.}
{Cây phượng ít hơn cây bằng lăng.}
{\True Cây phượng nhiều hơn cây bằng lăng là 4 cây ($12 - 8 = 4$).}
\loigiai{
\begin{itemize}
    \item Tổng: $15 + 12 + 8 = 35$ cây. Đúng.
    \item Bàng hơn bằng lăng: $15 - 8 = 7$ cây. Đúng.
    \item Phượng (12) nhiều hơn bằng lăng (8). Ý 3 Sai, Ý 4 Đúng.
\end{itemize}
\begin{itemize}
    \item a) \textbf{Đúng}.
    \item b) \textbf{Đúng}.
    \item c) \textbf{Sai}.
    \item d) \textbf{Đúng}.
\end{itemize}
}
\end{ex}

\textbf{Phần 3. Trả lời ngắn}
\setcounterpageref{ex}{0}

\begin{ex} %[3X15H73]
Lớp 3A có 18 bạn nam và 22 bạn nữ. Tính tổng số học sinh lớp 3A.
\shortans{40}
\loigiai{
\textbf{Đáp số:} $40$\\
$18 + 22 = 40$ học sinh.
}
\end{ex}

\begin{ex} %[3X15H73]
Một trang trại nuôi 150 con gà trống và 350 con gà mái. Hỏi đàn gà có tất cả bao nhiêu con?
\shortans{500}
\loigiai{
\textbf{Đáp số:} $500$\\
$150 + 350 = 500$ con.
}
\end{ex}

\begin{ex} %[3X15V73]
Thống kê số ki-lô-gam giấy vụn thu gom của 3 tổ: Tổ 1: $45\text{ kg}$, Tổ 2: $55\text{ kg}$, Tổ 3: $60\text{ kg}$. Tính tổng số ki-lô-gam giấy vụn cả 3 tổ thu gom.
\shortans{160}
\loigiai{
\textbf{Đáp số:} $160$\\
$45 + 55 + 60 = 160\text{ kg}$.
}
\end{ex}

\begin{ex} %[3X15V73]
Tổ 3 thu gom nhiều hơn tổ 1 bao nhiêu ki-lô-gam giấy vụn?
\shortans{15}
\loigiai{
\textbf{Đáp số:} $15$\\
$60 - 45 = 15\text{ kg}$.
}
\end{ex}

\begin{ex} %[3X15V73]
Một đội trồng cây: ngày 1 trồng 120 cây, ngày 2 trồng 150 cây, ngày 3 trồng 130 cây. Tính tổng số cây trồng trong 3 ngày.
\shortans{400}
\loigiai{
\textbf{Đáp số:} $400$\\
$120 + 150 + 130 = 400$ cây.
}
\end{ex}

\begin{ex} %[3X15C73]
Trong một giải chạy, số người về đích trong đợt 1 là 45 người, đợt 2 là 55 người, đợt 3 gấp đôi đợt 1. Tính tổng số người về đích trong 3 đợt.
\shortans{190}
\loigiai{
\textbf{Đáp số:} $190$\\
Đợt 3: $45 \times 2 = 90$ người.\\
Tổng: $45 + 55 + 90 = 190$ người.
}
\end{ex}
